% !TEX root = ../main.tex
\lecture{1}{2026-09-14}{Day 1 Intro}
\textbf{Vorlesungsfolien} \url{https://lms.uzh.ch/auth/RepositoryEntry/17903616206/CourseNode/114298086656257/path%3D~~Slides/0}
\section{Admin stuff}
\begin{itemize}
    \item Aufbau
    \subitem Corporate Finance (\textbf{INHALT DES 1. SEMESTERS})
        \subsubitem Bewertung von Investitionen
        \subsubitem Finanzierung von Unternehmen
        \subsubitem  Unternehmensbewertung
    \subitem Financial Econ
        \subsubitem Optimierung von Investitionsportfolios 
    \subitem Banking And Insurance (\textbf{INHALT DES 1. SEMESTERS})
        \subsubitem Bedeutung der Banken fuer die Ökonomie
    \subitem Quantitative Finance
    \subitem Sustainable Finance
    \item Vorlesung
        \subitem Werden aufgezeichned und Livestreamed 
    \item \textbf{Einzelaufgaben}
        \subitem 05.10 - 19.10.25 oder 26.10-09.11.26
    \item \textbf{Prüfung}
        \subitem 18.12.26 10 bis 11.30 Uhr
        \subitem 90 Minuten 
    \item Excel tutorial auf OLAT (IT-Support - Excel-Tutorial)
    \item Klicker AI Finance Chatbot mit Skript und Vorlesung 
\end{itemize}

\section{Was ist Finance?}
\textbf{Definitionen}
\begin{itemize}
    \item Verwaltung von Werten, Geld, Bankoperationen, Investitionen und Kredit
\end{itemize}
\break
\section{Bereiche von Finance}
\begin{itemize}
    \item VWL
        \subitem Zentralbank (Wertanlagen, Staatsschulden) 
    \item BWL
        \subitem Corporate Finance (Kredite, Aktien) 
    \item Mathematik
        \subitem Portfolio Management Quant   
\end{itemize}

\section{Intro Corporate Finance}
\textbf{Definition} Unternehmensfinanzierung umfasst alle Massnahmen, welche im Zusammenhang mit der Kap.beschaffung (Fremd, Eigenkap.), Kap.einsatz (interne Projekte), Kap.bewirtschaftung (Anlagen), und Rueckzahlung (Dividenden) getroffen werden müssen.
\subsection{Merger Baloise Helvetia}
Aus Kapitaleigentuemersicht 40\% Dividendenwachstum

\textbf{Fragen} Wie Resourcenallokation? Finanzierung von Investitionen, Groesse der Dividenden [Fluchtkapital]    


\textbf{Valuation, Merger, Aquisition als Bereich von Finance}
Bewertung des Kapitals und entsprechender Aktienpreis 

\subsection{Kursverläufe}
\textbf{SMI} Aktiensammlung der grössten schweizer Unternehmen

\subsection{Oberziele}
\begin{itemize}
    \item Grundbegriffe CF
    \item statistischen Investitionsverfahren
    \item Grunprinzip Time Value of Money, wann ist eine Geldeinheit mehr profitabel?
    \item Eigenschaften Fremd und Eigenkapital (Vor- nachteile)
    \item Finanzierungsformen 
    \item Unternehmensbewertungsverfahren
\end{itemize}

"Risikenschätzung, rückwirkend einfach aber in zukunft zB rechenzentren wenig risiko aber weiss nie wenn eine chinesische firma für 1/10 der Rechenleistung gleiche Resultate bringen kann"
[Solche Optimierung von Gebrauchswerten ist fuer den Kapitalist einen Albtraum, solange er/sie dessen Besitzer ist]


\section{Finanzwirtschaft. Zieldreieck}
Sicherheit (Langfristige ANlagerung oder Investition, risiko arm) - Liquiditaet (Mengegeld unmittelbar zur Verfügung) - Rentabilitaet (Profite)
Hauptaspekte, Tradeoffs $(\uparrow Sicherheit dann \downarrow Rentabilitaet)$

\subsection{Shareholder vs Stakeholder Value}
Unterscheidung, Priorisierung der Entlohnung Eigenkapital, Angestellten (Stakeholder) vs Shareholder (Fremdkapital)